% Generated draft: requires mathematical and source review.
\documentclass[11pt]{article}
\usepackage[margin=25mm]{geometry}
\usepackage{fontspec}
\setmainfont{Latin Modern Roman}
\usepackage{amsmath,amssymb,mathtools,unicode-math}
\setmathfont{Latin Modern Math}
\usepackage{enumitem,needspace,xurl,hyperref}
\setlength{\parindent}{0pt}
\setlength{\parskip}{5pt}
\setlength{\emergencystretch}{3em}
\newcommand{\N}{\mathbb{N}}
\newcommand{\Z}{\mathbb{Z}}
\newcommand{\Q}{\mathbb{Q}}
\newcommand{\R}{\mathbb{R}}
\newcommand{\C}{\mathbb{C}}
\newcommand{\F}{\mathbb{F}}
\newcommand{\A}{\mathbb{A}}
\newcommand{\PP}{\mathbb{P}}
\newcommand{\E}{\mathbb{E}}
\newcommand{\eps}{\varepsilon}
\newcommand{\dd}{\,\mathrm{d}}
\begin{document}
\section*{UnsolvedMath ID: 30002337}
\textbf{Source problem code:} catalogue entry
\subsection{Definitions and mathematical statement}
The supplied record collects six questions rather than specifying one complete conjecture. The definitions below describe standard interpretations and identify the additional data needed to make each question precise; these interpretations are not asserted to have been fixed by the source.

\begin{enumerate}
\item \textbf{Local minima of knot energies.} Let \(\mathcal X\) be a specified topological space of admissible embedded curves and let \(E:\mathcal X\to\R\cup\{+\infty\}\) be a specified energy. A finite-energy curve \(\gamma\in\mathcal X\) is a local minimum if there is a neighborhood \(U\) of \(\gamma\) such that
\[
E(\gamma)\leq E(\eta)\qquad\text{for every }\eta\in U.
\]
The source asks whether complexity methods can establish the existence of nontrivial local minima for knot energies. It does not specify the energy, admissible curves, topology, geometric normalization, meaning of complexity methods, or meaning of ``nontrivial.'' In particular, a minimum represented by a nontrivial knot and a local minimum that is not global are distinct possible interpretations. No particular existence assertion is determined by the supplied record.

\item \textbf{Connectedness of equilateral configuration spaces.} Under the usual spatial interpretation, an embedded equilateral arm with \(n\) edges of common length \(\ell>0\) is a vertex tuple \((v_0,\ldots,v_n)\in(\R^3)^{n+1}\) with \(\|v_i-v_{i-1}\|=\ell\) for \(1\leq i\leq n\), whose associated piecewise linear map from \([0,n]\) is injective. An embedded equilateral polygon with \(n\) edges is a cyclic vertex tuple with all cyclic edge lengths equal to \(\ell\), whose associated piecewise linear map from a circle is an embedding. Such a spatial polygon is unknotted if its image is ambient isotopic in \(\R^3\) to a round circle. For fixed parameters these tuples give configuration spaces with the subspace topology from their vertex coordinates, before any quotient or boundary conditions are imposed.

The source asks whether the spaces of embedded equilateral arms and unknotted equilateral polygons are connected. It does not explicitly fix ambient dimension, numbers of edges, endpoint constraints, identifications by rigid motions or vertex relabeling, or the quantifiers over these choices. The displayed spatial definitions therefore provide a possible setting, not a uniquely recovered formulation. The source says ``connected''; whether path connectedness is intended is also unresolved.

\item \textbf{Enumeration by stick number.} Under the usual interpretation for tame knots in \(\R^3\), a knot type is an ambient-isotopy class of embedded circles, and its stick number \(s(K)\) is the least number of straight edges in a polygonal representative. Once a convention for the class \(\mathcal K\) of knot types is fixed, the enumeration question is to determine
\[
N(m)=\#\{K\in\mathcal K:s(K)=m\},\qquad m\in\Z,\quad m\geq3.
\]
The source does not specify whether knots are oriented, whether mirror images are identified, or whether enumeration includes all knot types or a restricted class such as prime knots. These choices must remain unresolved.

\item \textbf{Limiting torsion laws.} An equilateral walk is a finite sequence of vertices with successive increments of a common length; an equilateral polygon additionally closes. Neither description alone specifies a probability distribution. The source asks for limiting torsion laws for confined equilateral walks and polygons, but supplies no confinement region, sampling measure, boundary treatment, definition of torsion, normalization, or limiting regime. Consequently no particular random variable or convergence-in-distribution statement can be specified from this record. Embeddedness is not explicitly required for these probabilistic objects.

\item \textbf{Mutation and arc index.} In the standard spatial setting, an arc presentation of a knot uses finitely many half-plane pages of an open book with a common binding axis, with the knot meeting each occupied page in one arc whose endpoints lie on the binding. The arc index \(\alpha(K)\) is the minimum number of occupied pages in such a presentation. The source asks whether mutation preserves arc index. If \(\mathcal M\) denotes the intended class of pairs of knot types related by the intended mutation operation, the question has the form
\[
\alpha(K)=\alpha(K')\qquad\text{for every }(K,K')\in\mathcal M.
\]
The record does not define mutation or its permitted variants. Ordinary Conway mutation is a possible intended interpretation, but cannot be selected conclusively from the supplied text. Any orientation or other restrictions on \(\mathcal M\) also require clarification.

\item \textbf{Frisch--Wasserman--Delbr\"uck exponent.} The source asks for the Frisch--Wasserman--Delbr\"uck exponent \(\varepsilon\). It gives no defining asymptotic relation, random-knot model, observable, or normalization. The name and symbol alone do not determine a numerical or asymptotic target; no formula defining \(\varepsilon\) is supplied here.
\end{enumerate}

The supplied status is \texttt{partially\_solved}. Its supplied dated assessment reports partial answers concerning random-polymer knotting and computational studies of equilateral knot spaces, while describing the general local-minimum, connectivity, enumeration, torsion, and arc-index questions as unresolved. This is source curation, not an independently verified mathematical conclusion.
\subsection{Short English statement}
This record groups six questions about knot-energy minima, connectedness of equilateral configuration spaces, counting knots by stick number, torsion distributions under confinement, mutation and arc index, and a Frisch--Wasserman--Delbr\"uck exponent. Several questions omit essential choices of spaces, constraints, probability models, or conventions. The formulations above preserve those gaps explicitly. The supplied partially solved label and literature assessment require independent review.
\subsection{Sources}
\begin{itemize}
\item Supplied local UnsolvedMath registry record: permanent ID \(30002337\), exact source problem code \texttt{catalogue entry}, title ``Geometric, Energetic, and Probabilistic Knot Problems,'' in the Oberwolfach Reports Open Problems collection. The record lists proposed year 2013 and supplies no named proposer.
\item Supplied report citation: \emph{Geometric Knot Theory (2014)}, Oberwolfach Reports, DOI 10.4171/owr/2013/22, pp.~1313--1358. \url{https://doi.org/10.4171/owr/2013/22}. The source describes Oberwolfach Reports as documenting mathematical workshops and their open-problem discussions. The relationship between the citation's 2014 label and the report's 2013 designation is not resolved by the supplied evidence; no individual problem page is identified.
\item Supplied literature source: \emph{On some knot energies involving Menger curvature}. \url{https://doi.org/10.1016/j.topol.2013.05.022}.
\item Supplied literature source: Diao--Pippenger--Sumners, \emph{On Random Knots}. \url{https://doi.org/10.1142/S0218216594000307}.
\item Supplied literature source: \emph{Oberwolfach Report 22/2013}. \url{https://doi.org/10.4171/owr/2013/22}.
\item Literature assessment date: \texttt{2026-08-21}. The partially solved label and the assessment of partial answers and remaining unresolved components are supplied dated curation derived from the supplied OWR curation, not independently verified here. The record explicitly requests independent verification of time-sensitive status and cited claims.
\end{itemize}
\end{document}
